\begin{table}[htbp]
\centering
\caption{Effect of NLW Bite on Log Wage Percentiles: Main DiD Results}
\label{tab:main}
\small
\begin{tabular}{lccccc}
\hline\hline
 & (1) & (2) & (3) & (4) & (5) \\
 & Log p10 & Log p25 & Log p50 & Log p60 & Log p90 \\
\hline
\textit{Dependent variable:} & & & & & \\
Bite $\times$ Post  & 0.1937*** & 0.2906*** & 0.2181*** & 0.1898*** & 0.1161* \\
 & (0.0174) & (0.0225) & (0.0295) & (0.0323) & (0.0620) \\
\hline
LA FE           & Yes & Yes & Yes & Yes & Yes \\
Year FE         & Yes & Yes & Yes & Yes & Yes \\
Clustered SE    & LA  & LA  & LA  & LA  & LA  \\
Observations    & 3,997 & 3,997 & 3,997 & 3,992 & 687 \\
$R^2$           & 0.980 & 0.961 & 0.939 & 0.933 & 0.973 \\
\hline\hline
\end{tabular}
\raggedright\footnotesize\textit{Notes:} All regressions include LA and year fixed effects. Standard errors clustered at the LA level in parentheses. Bite $\times$ Post is the interaction of the 2015 NMW bite ratio (£6.50/LA median wage) with an indicator for years 2016 and later (post-NLW introduction). Data: ASHE 2013--2023, all Local Authorities in England and Wales. *** $p<0.01$, ** $p<0.05$, * $p<0.10$.
\end{table}
